\lecture{1}{Tue 09 Jan 2024 10:34}{Linear Algebra in Hilbert Spaces}

Reading: all front matter in chapter 1, asap. This is for motivation and quick intro to background of QC. The mathematical content is in Chap. 2.1.

\textbf{Why QC?}

\subsubsection{For Computer Scientists}


\begin{itemize}
	\item For CS: QC threatens the validity of the \textit{extended Church-Turing thesis}.
		The CTT says that anything that is computable can be computed by a Turing machine.
		The Extended CTT says that everything that is \textbf{efficiently} computable is \textbf{efficiently} computable using a Turing machine. \textit{Efficient} means the time complexity is polynomial.
	
		QC threatens this from various angles. Most famously, via the \textit{Shor's algorithm}.

		There is an efficient probabilistic algorithm in the 90's for determining primality; two years later we got an efficient deterministic algorithm. It is suspected that BQP = P, so every probabilistic argument can be determinized. 

	The Shor's algorithm says: For any integer $N$ expressed in binary (so input of $\left\lceil \log\left( N \right)  \right\rceil $ , there is a polynomial time algorithm to factor $N$ on a quantum computer.

	Currently the argument is probabilistic; and 

	BQP are not expected to overlap NP-complete, but some NP problems are contained in BQP. Proving the extended CTT false would require some breakthrough in computational complexity theory.

	Another CS reason to care about QC is cryptography.


\end{itemize}

\subsubsection{For Physicists}

\begin{itemize}
	\item A computational perspective on quantum mechanics is very fruitful, in condensed matter physics, HEP, quantum gravity.
\end{itemize}

\subsubsection{For Mathematicians}
\begin{itemize}
	\item See CS answer
\end{itemize}

\subsubsection{For Engineers}
\begin{itemize}
	\item A physical perspective on computation is fruitful. 
\end{itemize}

I will try to justify all these. Moreover, I want to give a sense of \textbf{what QC cannot do}.

\subsection{A nugget of perspective}

There is a philosophy in mathematics called \textit{quantum mathematics}: Quantum Mechanics is a set of rules that should be understood as a rigorous mathematical theory of ``quantum probability''. 

There exist many types of QM systems in nature, but they have some common nature that is captured by the \textbf{axioms of quantum mechanics.}

(This is established in 1920s). 

Why are these axioms? Take a physics QM class. They will explain the experiments that show nature is inconsistent with classical mechanics.

\subsubsection{Analogy between classical probability and QM}
\begin{itemize}
	\item Sample space $\Omega$ $\leftrightarrow$ state space (finite dimensional Hilbert space) $\mathscr{H}$ 
	\item probability $0 \le  p \le  1$ $\leftrightarrow$ amplitude $z \in \mathbb{C}, |z| < 1$
	\item distribution on $\Omega$ $\leftrightarrow$ A non-zero vector in $\mathscr{H}$.
	\item random variable $X: \Omega \to S$ $\leftrightarrow$ \textit{observable} (Hermitian linear operator $A: \mathscr{H} \to \mathscr{H}$ )
	\item outcome of a random variable: $s \in S$ $\leftrightarrow$ outcome of observation: eigenvalue of $A$
	\item symmetry of  $\Omega$ is a bijection $\leftrightarrow$ symmetry of $\mathscr{H}$ is a \textit{unitary transformation}
	\item Two-part system $\Omega_1 \cup \Omega_2$ $\leftrightarrow$ $\mathscr{H}_1 \otimes \mathscr{H}_2$
\end{itemize}

\section{Linear Algebra Review}

The most important mathematical structure in QM is a \textit{Hilbert Space}.

Let's recall

\begin{definition}
A vector space $\mathscr{H}$ over $\mathbb{C}$ is a set $\mathscr{H}$ with two additional structure:
\begin{itemize}
	\item  $+: \mathscr{H} \times \mathscr{H} \to  \mathscr{H}$
	\item $\cdot : \mathbb{C} \times  \mathscr{H} \to  \mathscr{H}$
\end{itemize}
Additional properties:
\begin{itemize}
	\item (A1) $+$ is associative
	\item (A2) exists (unique)  $0 \in  \mathscr{H}$ that is additive identity (Monoid)
	\item (A3) there exists additive inverses. (Group)
	\item (A4) $+$ is commutative (Abelian Group)
	\item (V1) there is a multiplicative identity $1 \in \mathbb{C}$.
	\item (V2) $\cdot$ is associative.
	\item (V3) scalar addition distribute over multiplication
	\item (V4) vector addition distribute over multiplication
\end{itemize}
\end{definition}

\begin{definition}
	If $\mathscr{H}$ is a $\mathbb{C}$-vector space, then an inner product is a function
	\[
		\left\langle -, - \right\rangle : \mathscr{H} \times \mathscr{H} \to \mathbb{C}
	\] 
	that satisfies
	\begin{enumerate}
		\item linear in second variable: $\left\langle a, z_1 b + z_2 c \right\rangle = z_1 \left\langle a, b \right\rangle  + z_2 \left\langle a, c \right\rangle $
		\item conjugate symmetry: $\left\langle b, a \right\rangle  = \left\langle a, b \right\rangle^{*}$
		\item positive definite: $\left\langle a, a \right\rangle \ge 0$. Also, $\left\langle a, a \right\rangle =0 \iff  a =0$.
	\end{enumerate}
\end{definition}
The first two properties give us algebraic properties; the third property allows us to do analysis. 

We define the \textit{length }of a vector $a \in \mathscr{H}$ is
\[
	\|a\| = \sqrt{\left\langle a, a \right\rangle } 
.\] 
The \textit{distance} between $a$ and $b$ is $\|b - a\|$.

\begin{theorem}[Cauchy-Schwatz Inequality]
	\label{lem:CauchySchwatzInequality}
	\[
		\left| \left\langle a, b \right\rangle  \right| ^2 \le 
		\|a\|^2 \|b\|^2
	.\] 
\end{theorem}

\begin{corollary}[Triangle Inequality]
	\[
		\|a + b\| \le  \|a\| + \|b\|
	.\] 
	Thus, \textit{every Hilbert space is a metric space.}
\end{corollary}

\begin{definition}[Hilbert space]
	\label{defn:HilbertSpace}
	
	A \textit{Hilbert} space is a complex inner product space such that the metric is complete.
\end{definition}

\begin{fact}
	All finite dimensional complex inner product spaces are automatically complete, and hence Hilbert.
\end{fact}

\begin{definition}[Adjoint]
	\label{defn:adjoint}
Let $A: \mathscr{H} \to \mathscr{H}$ be a linear transformation. Define the \textit{adjoint} of $A$ is the (unique! see H) linear transformation $A^*: \mathscr{H} \to  \mathscr{H}$ defined by the property
\[
	\left\langle A^* a, b \right\rangle  = \left\langle a, Ab \right\rangle 
.\] 
\end{definition}

\begin{lemma}
	$A^*$ is well-defined and $(AB)^* = B^* A^*$
\end{lemma}

\begin{definition}
	We say $A$ is 
	\begin{enumerate}
		\item \textit{Hermitian} (or self-adjoint) if $A^* = A$
		\item \textit{Unitary} if $A^*  = A^{-1}$
		\item \textit{Normal} if $A A^* = A^* A$
	\end{enumerate}
\end{definition}
(Check) Hermitian or Unitary implies Normal.

\todo{Need to play with those types of operators and gain intuition.}

\begin{lemma}
	\label{lem:UnitaryCondition}
	$U: \mathscr{H} \to  \mathscr{H}$ is unitary iff
	\[
		\left\langle Ux, Uy \right\rangle  = \left\langle x, y \right\rangle  
		\qquad \forall x, y \in \mathscr{H}
	.\] 
\end{lemma}

\begin{definition}[Unitarily Diagonalizable]
	\label{defn:UnitarilyDiagonalizable}
	We say $A: \mathscr{H} \to  \mathscr{H}$ is \textit{unitarily diagonalizable} if there exists an orthonormal basis, such that the matrix of $A$ is diagonal in that basis.
\end{definition}

\begin{theorem}[Spectral Theorem in Finite Dimensions]
	\label{thm:SpectralThmFiniteDim}
	Let $\mathscr{H}$ be a finite dimensional Hilbert space. Then an operator $A: \mathscr{H} \to  \mathscr{H}$ is normal iff $A$ is unitarily diagonalizable.
\end{theorem}

\begin{corollary}
	\begin{enumerate}
		\item Hermitian iff unitarily diagonalizable with real eigenvalues
		\item Unitary if unitarily diagonalizable with all eigenvalues unitary (length $1$ )
	\end{enumerate}
\end{corollary}

Next lecture: we talk about tensor products.
